% Pollis drill-down: Explore with Pollis AI+, the chat model sent out into the world with tools and skills, beside
% Ask Pollis AI+, which answers from Pollis's own documents; and how each harness is measured, model by model, so that
% the progress of the harnesses can be seen (4 November 2026).
% Written before the code, as these are to work: the plan to build them from. Check every statement against Pollis
% once it is built.
% Compile with: latexmk -pdf workflow_explore_ai.tex
% Shared look for the one-page Pollis workflow handouts (same fonts and colours as the brochure).
\documentclass[10pt,a4paper]{article}

\usepackage[a4paper,left=17mm,right=17mm,top=14mm,bottom=16mm,footskip=8mm]{geometry}
\usepackage[T1]{fontenc}
\usepackage[default,semibold]{sourcesanspro}
\usepackage{microtype}
\usepackage{graphicx}
\usepackage[table]{xcolor}
\usepackage{tikz}
\usepackage[most]{tcolorbox}
\usepackage{fontawesome5}
\usepackage{enumitem}
\usepackage{fancyhdr}
\usepackage[hidelinks]{hyperref}

\graphicspath{{images/}}

% ---------------------------------------------------------------- colours (as in the brochure)
\definecolor{navy}{HTML}{1C3A5E}
\definecolor{pblue}{HTML}{3B78B5}
\definecolor{amber}{HTML}{F2B33D}
\definecolor{deepamber}{HTML}{D9902A}
\definecolor{pale}{HTML}{EEF3F9}
\definecolor{ink}{HTML}{2A2F36}
\definecolor{logonavy}{HTML}{173D63}
\definecolor{logodark}{HTML}{C98A1E}
\definecolor{logoamber}{HTML}{F2B446}
\definecolor{logocream}{HTML}{FEEFD2}
\definecolor{soft}{HTML}{5E6773}
\color{ink}

\setlength{\parindent}{0pt}
\setlength{\parskip}{3pt}
\setlist{nosep,leftmargin=*,itemsep=2pt}

\pagestyle{fancy}
\fancyhf{}
\renewcommand{\headrulewidth}{0pt}
\fancyfoot[L]{\footnotesize\color{soft}\textbf{Pollis}\enspace Bright Insights through Exploration, Modelling, Simulation and Optimisation}
\fancyfoot[R]{\footnotesize\color{soft}github.com/saragga/pollis/releases/latest}

% The Pollis logo, redrawn so it prints sharply at any size
\newcommand{\pollislogo}[1]{%
  \begin{tikzpicture}[x=#1/512,y=-#1/512,baseline=(current bounding box.center)]
    \fill[logonavy] (0,0) rectangle (512,512);
    \fill[logodark] (36,422) -- (104,256) -- (172,422) -- cycle;
    \fill[logoamber] (104,422) -- (203,180) -- (302,422) -- cycle;
    \fill[logocream] (203,422) -- (339,89) -- (476,422) -- cycle;
  \end{tikzpicture}}

% Top of the page: logo, name, kind of document (Workflow unless given, e.g. [Drill-Down]) and its subject;
% then the question and a one-line promise
\newcommand{\workflowhead}[4][Workflow]{%
  \thispagestyle{fancy}%
  \begin{minipage}[c]{13mm}\pollislogo{11mm}\end{minipage}%
  \begin{minipage}[c]{\dimexpr\linewidth-13mm\relax}
    {\color{navy}\Large\bfseries Pollis}\enspace{\color{soft}#1}\hfill
    {\color{deepamber}\small\bfseries\MakeUppercase{#2}}
  \end{minipage}\par\vspace{7mm}
  {\color{navy}\fontsize{22}{26}\selectfont\bfseries #3}\par\vspace{4pt}
  {\color{soft}\large #4}\par\vspace{5mm}}

% A menu path: \menu{Explore}{Access Data Libraries}... is written as \path{Explore \gt Access Data Libraries}
\newcommand{\gt}{\,{\color{deepamber}\bfseries ›}\,}
\newcommand{\menupath}[1]{{\footnotesize\bfseries\color{navy}#1}}

% A screenshot if images/<file> exists, otherwise an orange placeholder that says what to capture
\newlength{\shotheight}\setlength{\shotheight}{43mm}
\newcommand{\stepshot}[2]{%
  \IfFileExists{images/#1}{%
    \tcbox[enhanced, colback=black, colframe=black!70, boxrule=0.4pt, arc=2pt,
      left=0pt,right=0pt,top=0pt,bottom=0pt, boxsep=0pt, halign=center]{%
      \includegraphics[width=\dimexpr\linewidth-1pt\relax,height=\shotheight,keepaspectratio]{#1}}}{%
    \begin{tcolorbox}[enhanced, colback=white, colframe=orange!80, boxrule=0.6pt, arc=2pt,
      height=\shotheight, valign=center, halign=center, left=6pt, right=6pt]
      {\color{orange}\Large\faCamera}\par\smallskip
      {\color{orange}\footnotesize\bfseries Screenshot: #2}\par\smallskip
      {\color{orange!80!black}\scriptsize images/#1}
    \end{tcolorbox}}}

% One step: number, title, menu path, image file, what to capture, what it means
\newcommand{\step}[6]{%
  \begin{tcolorbox}
    {\color{deepamber}\Large\bfseries #1}\enspace{\color{navy}\large\bfseries #2}\par\vspace{1pt}
    \menupath{#3}\par\vspace{4pt}
    \stepshot{#4}{#5}\par\vspace{3pt}
    {\small #6}
  \end{tcolorbox}}

\newenvironment{steps}{%
  \begin{tcbraster}[raster columns=2, raster equal height=rows, raster column skip=4mm,
    raster row skip=4mm, enhanced, colback=pale, colframe=pale, boxrule=0pt, sharp corners,
    left=6pt,right=6pt,top=6pt,bottom=6pt, fontupper=\raggedright]}{\end{tcbraster}}

% Boxes for the bottom of the page
\newtcolorbox{pbox}[2][]{enhanced, colback=pale, colframe=amber, boxrule=0pt,
  leftrule=3pt, sharp corners, left=7pt, right=7pt, top=6pt, bottom=6pt,
  fontupper=\small\raggedright, before upper={{\color{navy}\normalsize\bfseries #2}\par\smallskip}, #1}
\newtcolorbox{nbox}[1][]{enhanced, colback=navy, colframe=navy, boxrule=0pt,
  sharp corners, left=8pt, right=8pt, top=6pt, bottom=6pt, coltext=white,
  fontupper=\small\raggedright, #1}

% A Next Steps item inside the bottom boxes
\newcommand{\nsitem}[2]{\item[\color{deepamber}\faChevronRight] \textbf{#1}: #2}

% The page of each heading, for Ask Pollis
% Pollis: the page of each heading, for Ask Pollis. A document that inputs this file, at the end of its preamble, writes
% <jobname>.labels beside its PDF, a line page|level|heading for each heading, the page counted from the first, the
% cover, as 1. compile_documents.jl copies it to ~/.pollis/documents, and ask_index.jl cuts the document into passages
% under its headings with it. \asklabel{1}{Title} marks a topic or a page, \asklabel{2}{Title} a heading inside it; the
% sections and subsections of a document that has them are marked as they go into its contents. A package vignette
% inputs it before its \documentclass (compile_vignette.jl), so the contents are taken over when the document begins,
% after the packages have taken them over.
\makeatletter
\newwrite\pollis@labels
\immediate\openout\pollis@labels=\jobname.labels
\newcommand{\asklabel}[2]{\write\pollis@labels{\the\ReadonlyShipoutCounter|#1|\unexpanded{#2}}}
\def\pollis@toc{toc}
\newcommand{\pollis@contentsline}[3]{%
  \def\pollis@file{#1}%
  \ifx\pollis@file\pollis@toc\asklabel{#2}{#3}\fi
  \pollis@addcontentsline{#1}{#2}{#3}}
% The headings of the New Methods style of the vignettes, \nmsection and \nmsub, which go into no contents
\def\pollis@asksection#1{\pollis@nmsection{#1}\asklabel{1}{#1}}
\def\pollis@asksub#1{\pollis@nmsub{#1}\asklabel{2}{#1}}
\AddToHook{begindocument/end}{%
  \let\pollis@addcontentsline\addcontentsline\let\addcontentsline\pollis@contentsline
  \ifdefined\nmsection\let\pollis@nmsection\nmsection\let\nmsection\pollis@asksection\fi
  \ifdefined\nmsub\let\pollis@nmsub\nmsub\let\nmsub\pollis@asksub\fi}
\makeatother



\usepackage{titlesec}
\usepackage{tabularx}
\usepackage{booktabs}
\usepackage{listings}

\fancyhf{}
\fancyfoot[L]{\footnotesize\color{soft}\textbf{Pollis}\enspace Drill-down: Explore with Pollis AI+}
\fancyfoot[R]{\footnotesize\color{soft}\thepage}

\titleformat{\section}{\color{navy}\large\bfseries}{\color{deepamber}\thesection}{0.6em}{}
\titleformat{\subsection}{\color{navy}\normalsize\bfseries}{\color{deepamber}\thesubsection}{0.5em}{}
\titlespacing*{\section}{0pt}{12pt}{4pt}
\titlespacing*{\subsection}{0pt}{8pt}{2pt}
\setlength{\parskip}{4pt}

\lstset{basicstyle=\ttfamily\footnotesize, columns=fullflexible, keepspaces=true, breaklines=true,
  backgroundcolor=\color{pale}, frame=leftline, rulecolor=\color{amber}, framerule=2pt, xleftmargin=6pt,
  framexleftmargin=6pt, aboveskip=4pt, belowskip=4pt, upquote=true, tabsize=2}

\usepackage{xurl}
\newcommand{\file}[1]{{\small\nolinkurl{#1}}}
\newcolumntype{L}{>{\raggedright\arraybackslash}X}

\begin{document}

\workflowhead[Drill-Down]{Explore with Pollis AI+}
  {The same small model, sent out into the world with tools and skills}
  {Two harnesses for one chat model: Ask Pollis AI+, which answers from Pollis's own documents, and Explore with
   Pollis AI+, which searches, reads and computes beyond them; the skills that make a small model on this computer
   reliable at both; and how each harness is measured, model by model, so that its progress can be seen.}

\begin{nbox}
{\color{amber}\bfseries Drill-down.}\enspace
The plan of Explore with Pollis AI+ and of the measure of its harnesses, written before the code. It builds on the
drill-down of Ask Pollis, whose Evaluate Answering Models it extends (\S\ref{sec:measure}).
\end{nbox}

\tableofcontents

\section{Two harnesses, one model}\label{sec:two}

A chat model alone answers from what it happens to know. A \textbf{harness} is everything Pollis puts around it: the
instruction, what it is given to read, the tools it may call, the steps it is led through, and the checks on what
it writes. The model can be the same for two jobs; the harness is what fits it to each.

\begin{center}\small
\renewcommand{\arraystretch}{1.25}
\begin{tabularx}{\linewidth}{@{}p{0.2\linewidth}LL@{}}
\toprule
& \textbf{Ask Pollis AI+} & \textbf{Explore with Pollis AI+}\\
\midrule
Answers about & Pollis, its toolboxes, its packages and Julia & the world beyond Pollis: papers, data, the web, your
  own files\\
Reads & a closed set of documents, indexed once (the drill-down of Ask Pollis) & what its tools bring back for the
  question: search results, pages, PDFs, tables\\
Harness & retrieval: the eight passages found, and an instruction to keep to them and cite them & skills: fixed
  steps, each with its tools, and checks between them (\S\ref{sec:skills})\\
Tools & none & web search and fetch, scholarly search, the PDF Viewer and its Poppler tools, Julia, your references,
  the MCP servers you set up (\S\ref{sec:tools})\\
Menu & \menupath{Help\gt Ask Pollis} & \menupath{Explore\gt Explore with Pollis AI+}\\
Sends out & nothing with a local model & with a local model, only the searches and the pages fetched; with a cloud
  model, also what it reads\\
\bottomrule
\end{tabularx}
\end{center}

Ask Pollis AI+ is help: it answers what Pollis is and how to use it. Explore with Pollis AI+ is exploration, which is
why it goes in the Explore menu, with the Browser at its top.

\section{A harness for a small model}\label{sec:small}

A model with two or four billion weights is good at small, well-posed tasks: write a search query, choose the
passage that answers, take the numbers out of a table, summarise a section. It is poor at long free loops, where it
must choose among many tools and plan several steps ahead, since a mistake at one step is carried into every step
after it. So the harness does the planning, and the model the small decisions.

\begin{center}\small
\renewcommand{\arraystretch}{1.25}
\begin{tabularx}{\linewidth}{@{}p{0.27\linewidth}p{0.2\linewidth}L@{}}
\toprule
\textbf{Model} & \textbf{Harness} & \textbf{The model decides}\\
\midrule
About 2B weights (effective), such as Gemma E2B & fixed workflow & only what each step asks, in a form fixed in advance\\
About 4B (effective), such as Gemma E4B & fixed workflow & the same, with longer steps: a whole section to summarise\\
About 12B, such as Gemma 12B & guided & which of two or three tools a step calls, and when a step is done\\
A cloud model, such as Claude & free & the steps themselves, with every tool of the skill\\
\bottomrule
\end{tabularx}
\end{center}

\textbf{Effective} weights: E2B and E4B hold more weights in memory than their names say, which the memory figure of
\S\ref{sec:figures} shows.

\subsection{Answers in a fixed form}
At each step of a workflow, the model's answer is held to a JSON schema: llama-server, which runs the local models,
then lets the model write only what fits the schema. A small model asked for \texttt{\{"query": "..."\}} cannot answer
with a paragraph instead, so the step that reads its answer never fails on its form; only its content can be wrong,
and the next step's check is for that.

\subsection{Checks between the steps}
Each step's output is checked before the next step uses it, by the harness, not by the model: a query that is empty,
a page that has no text, a number that is not in the table it was taken from, a citation to a source that was not
read. A step that fails its check is run once more with what was wrong said in its instruction, and then stops the
skill with what it has, saying which step failed.

\section{The tools}\label{sec:tools}

\begin{center}\small
\renewcommand{\arraystretch}{1.25}
\begin{tabularx}{\linewidth}{@{}p{0.24\linewidth}L@{}}
\toprule
\textbf{Tool} & \textbf{What it does}\\
\midrule
Web search & a query in, a list of titles, links and snippets out\\
Fetch & a link in, the text of the page out, without its menus and adverts\\
Scholarly search & Semantic Scholar, OpenAlex and arXiv: works, with their authors, abstracts, citations and DOIs\\
Read a PDF & the text of a PDF, page by page, with Poppler's \file{pdftotext}, as the PDF Viewer's Extract Text gives it\\
Julia & a computation whose result has to be right: a number in a summary, the solution of a problem\\
Your references & the works of PollisReferences and your library: is a work there, and with what key\\
MCP servers & those of \menupath{Compose\gt MCP Servers}, when a skill names them\\
LaTeX & a document compiled by PollisTeX, for the skills that give a PDF\\
\bottomrule
\end{tabularx}
\end{center}

A tool is called by the harness, at the step that needs it, for the small models, and by the model for the cloud one.
Every call is shown in the answer, with what went in and what came out, so that you can see what was read.

\section{The skills}\label{sec:skills}

A skill is a folder in the format of Agent Skills: a \file{SKILL.md} that says what the skill is for and its steps,
with the files those steps need. Pollis already uses this format for its agent skills (\menupath{Compose\gt Agent
Skills}), so the skills of Explore with Pollis AI+ can be read by a coding agent too. A skill adds to it, in its
front matter, the steps of the fixed workflow, each with its tool, its schema and its check, for the small models.

\begin{lstlisting}
---
name: summarise-a-paper
description: The question, data, method, results and limits of a paper, with the pages they come from, as a PDF.
steps:
  - read:      { tool: read-pdf }
  - outline:   { schema: sections.json,  check: every section has pages }
  - summarise: { schema: summary.json,   check: every page cited exists, every number is in the text }
  - write:     { tool: latex, template: summary.tex }
---
\end{lstlisting}

\subsection{The first skills}
The first skills are those of the Self-Study toolbox, in the order its Guide gives, and one of Explore:
\begin{center}\small
\renewcommand{\arraystretch}{1.25}
\begin{tabularx}{\linewidth}{@{}p{0.26\linewidth}L@{}}
\toprule
\textbf{Skill} & \textbf{Steps}\\
\midrule
Make More Problems & read the problems given; name the skill each tests; write new ones at the same level; solve
  each, in Julia when it is a computation; keep those whose solution checks; write the PDF\\
Summarise a Paper & read the PDF; find its sections; summarise question, data, method, results and limits, each with
  its pages; check every number against the text; write the PDF\\
Explain a Topic in Depth & take the paragraph or the area; list its concepts; explain each, with the steps of its
  argument and an example; write the PDF\\
Text into Slides & read the text; find its main points; one slide each, with at most five lines; compile with
  PollisTeX; keep the LaTeX source\\
Slides into Prose and Audio & read the slides; write the text that joins them, slide by slide; read it aloud with a
  voice model; write the PDF and the audio file\\
Find Papers on a Question & write two or three queries; search the scholarly sources; keep the works that answer;
  summarise each in two lines, with its DOI; mark those already in your references\\
\bottomrule
\end{tabularx}
\end{center}

\section{Measuring the harnesses}\label{sec:measure}

Evaluate Answering Models compares models on the same questions, the same passages and the same instruction, so that
what is compared is the model alone. To see what a harness adds, the comparison turns around: the model is held, and
the harness changes. The open point \emph{With and without the documents}, of the drill-down of Ask Pollis, is the
first case of it: a model with no harness, and with the passages.

\subsection{The harness as a dimension}
A run chooses models, as now, and harnesses: \textbf{None}, the model alone, and the versions of the harness of the
test set, each with its number. A harness version is frozen once it is used in a run: its instruction, its steps,
its schemas and its checks are kept with the run, so that a run months later can be set beside it. Each row of the
figures is then a model with a harness:

\begin{center}\small
\renewcommand{\arraystretch}{1.25}
\begin{tabular}{@{}llcccc@{}}
\toprule
\textbf{Model} & \textbf{Harness} & \textbf{Right} & \textbf{Seconds} & \textbf{Tokens} & \textbf{Memory}\\
\midrule
Gemma E2B & None & & & & \\
Gemma E2B & Version 1 & & & & \\
Gemma E2B & Version 2 & & & & \\
Gemma E4B & None, 1, 2 & & & & \\
Gemma 12B & None, 1, 2 & & & & \\
Claude & None, 2 & & & & \\
\bottomrule
\end{tabular}
\end{center}

The cloud model is run with the harness too: that shows how much each model gains from it, which is the figure that
matters, since the harnesses are what Pollis makes.

\subsection{The figures}\label{sec:figures}
Beside those Evaluate Answering Models gives now, the score from 1 to 5, its five parts, the seconds and the tokens:
\begin{center}\small
\renewcommand{\arraystretch}{1.25}
\begin{tabularx}{\linewidth}{@{}p{0.22\linewidth}L@{}}
\toprule
\textbf{Figure} & \textbf{What it says}\\
\midrule
Right & the share of answers the judge finds agree with the right answer of the test set; for a skill, the share
  whose output meets its rubric (\S\ref{sec:sets})\\
Gain & for a harness, Right minus that of the same model with None\\
Steps, Tool Calls & per answer, how many steps were run and tools called, and how many failed their check\\
Tokens & per answer, read and written, the tools' results included, since they are read too\\
Memory & for a local model, the most memory it used, which for E2B and E4B is more than their names say\\
Machine & the computer a run was made on, its chip and its memory, kept with the run\\
\bottomrule
\end{tabularx}
\end{center}

\subsection{Comparable over months}
\begin{itemize}
  \item \textbf{A frozen test set}, and a new one under a new name, never a changed one: Standard stays as it is.
  \item \textbf{A frozen judge}: the judge of a series of runs stays the same model, named with its version; a new
    judge starts a new series.
  \item \textbf{The machine}: local seconds are compared only between runs on the same machine.
  \item \textbf{Frozen sources for Explore}: the web changes its answers, so a test set of Explore keeps the pages
    and PDFs its questions were made from; a run reads those, not the web, unless you choose \textbf{Live}.
\end{itemize}

\subsection{The test sets}\label{sec:sets}
\begin{itemize}
  \item \textbf{Ask Pollis}: Standard, the 100 questions Pollis comes with, as now.
  \item \textbf{Explore}: questions whose answers are in open-access papers and pages kept with the test set: a
    number, a method, a finding, an author.
  \item \textbf{Self-Study}: tasks with a rubric rather than one right answer. A summary is right when it has the
    paper's question, data, method, results and limits, each with a page that holds it; a problem when its solution
    checks, in Julia for a computation; slides when each main point of the text is on one. The Self-Study toolbox
    brings this test set, and its Guide shows the table of \S\ref{sec:measure} for its tasks.
\end{itemize}

\subsection{What to expect}
On closed tasks, where the answer is in what the model is given, a summary of a paper you give it or a question about
Pollis, a well-harnessed model of four or twelve billion weights can come close to a cloud model: the harness finds
and checks, and the model only has to read. On open reasoning, where nothing given holds the answer, the gap stays
wide, and grows as the cloud models improve. The figures show where a model on this computer is already enough, and
where it is not yet, and each new harness version moves that line.

\section{The steps}\label{sec:steps}
\begin{enumerate}[label=\arabic*.]
  \item \textbf{The harness as a dimension} of Evaluate Answering Models: None and the harness of Ask Pollis, with
    the figures Right and Gain, on Standard. It measures what is there today, before anything new.
  \item \textbf{The skill runner}: a skill's steps, its schemas through llama-server, its checks, its tools called
    by the harness; shown step by step in the answer.
  \item \textbf{Make More Problems}, the first skill, and the Self-Study test set with its rubric.
  \item \textbf{Read a PDF, Julia and LaTeX} as tools, and \textbf{Summarise a Paper}.
  \item \textbf{Explore with Pollis AI+} in the Explore menu, with \textbf{Explain a Topic in Depth} and
    \textbf{Text into Slides}.
  \item \textbf{Web search, fetch and scholarly search}, with \textbf{Find Papers on a Question} and the Explore
    test set, with its frozen sources.
  \item \textbf{A voice model}, with \textbf{Slides into Prose and Audio}.
\end{enumerate}

\section{Open points}\label{sec:open}
\begin{itemize}
  \item \textbf{Web search}: a service with a key of yours, a free one with limits, or a search engine run on this
    computer; and what the hover of the answer says is sent out.
  \item \textbf{Tool calls of the 12B model}: whether it chooses among two or three tools reliably enough to be
    guided rather than led, which the figures of \S\ref{sec:figures} answer.
  \item \textbf{Memory}: a 12B model and a page or a PDF in its context may not fit a laptop with 16 GB; the warning
    of Evaluate Answering Models covers runs, and the skill runner needs one of its own.
  \item \textbf{The voice model}, for Read Aloud and for the audio of Slides into Prose and Audio: one made for speech
    or a general one, and which is the default.
  \item \textbf{The judge of a rubric}: a model reading a summary against a rubric judges less surely than one
    comparing two answers; a second judge (as in Evaluate Answering Models) and checks done by the harness, such as the
    pages cited, make it surer.
\end{itemize}

\end{document}
